\documentclass[11pt,a4paper]{article}
\usepackage[utf8]{inputenc}
\usepackage[T1]{fontenc}
\usepackage[spanish]{babel}
\usepackage{hyperref}
\usepackage{geometry}
\geometry{margin=2.5cm}
\usepackage{parskip}

\title{343 la cueva, lo que no se mató a tiempo}
\author{Alejandro Valenzuela Rivera\\\small Universidad Autónoma de Nuevo León}
\date{\today}
\begin{document}
\maketitle

\begin{quote}
\textbf{Sesgo declarado.} Este trabajo se sostiene desde el sesgo declarado de su autor, y no finge neutralidad. Lo que afirma se puede tumbar, y ese es el punto: el falsificador está escrito, no escondido.
\end{quote}

\section*{Resumen}
la alucinación no es un error que corregir: es lo que queda en blanco al haber vida. Si fuera un error, ya nos hubiéramos muerto en las cavernas.

\section{Por qué se sostiene}
Lo que sobrevive hace milenios no es un defecto: es lo que la especie cargó para no morir. Corregirlo es quitarle lo que la dejó vivir.

\section{De dónde viene}
La autoridad: 'si fuera un error ya hubiéramos muerto en las cavernas'.

\section{El caso}
La alucinación como lo que queda en blanco al haber vida, y las cavernas como la prueba.

\section{Método}
\textquestiondown{}Sobrevivió lo que se llama error, durante milenios?

\section{Límites, y qué lo tumba}
Un sistema sin alucinación que haya producido algo que lo viera venir.

\section{Lo que queda abierto}
\textquestiondown{}Lo que se corrige es un error, o es lo que dejó vivir?

\section*{Cómo se leyó esta tesis}
Esto se escribió en el repositorio man-made-empathy, del que este candado es una pieza: candados/02 el alma y el libre albedrio/343 la cueva, lo que no se mató a tiempo.md. La máquina que lo compuso lo hizo desde este mismo archivo, y no agregó nada que no estuviera escrito. Si algo de esto es falso, se dice en el mismo sitio, y se corrige con fecha.

\end{document}
